% Source: https://artofproblemsolving.com/community/c3h2865423_aint_need_calculus
% Stable identity: aops-topic-2865423-q01
For $0<x<\dfrac\pi2$, in the unit circle below, triangle $OAP$ lies inside sector $OAP$, which lies inside triangle $OAT$ (where $AT$ is tangent at $A$).
\begin{asy}
usepackage("amsmath");
size(4cm);
defaultpen(fontsize(10pt));
real t=40;pair O=(0,0),A=(1,0),P=dir(t),Q=(1,tan(radians(t)));
draw(arc(O,1,-5,70));draw(O--A--Q--cycle);draw(O--P);draw(P--(P.x,0),dashed);
draw(arc(O,0.25,0,t));label("$x$",0.25*dir(t/2),E);
label("$O$",O,SW);label("$A$",A,SE);label("$P$",P,NW);label("$T$",Q,NE);
label("$\sin x$",(P.x,P.y/2),W);label("$\tan x$",(1,Q.y/2),E);
\end{asy}

Comparing areas, $\dfrac12\sin x<\dfrac12x<\dfrac12\tan x$. So $\cos x<\dfrac{\sin x}{x}<1$.

Put $x=\dfrac\pi n$ for $n\ge3$ and multiply by $\pi$:
\[\pi\cos\dfrac\pi n<n\sin\dfrac\pi n<\pi.\]
As $n\to\infty$, $\cos\dfrac\pi n\to1$, so $k=\pi$. Then $100k=100\pi=314.159\ldots$, and the answer is $\boxed{314}$.
